\section{Equality in four spatial dimensions}
\label{q4:sec:setup}

Equality concerns the initial hypersurface, including its second fundamental
form. A limiting scalar-flat comparison metric does not by itself identify
that hypersurface. We instead apply the numerical inequality to nearby
original data. The resulting variational identity produces a causal
stationary field on the original exterior. We remove its twist, classify
the complete static base, and then recover the original hypersurface in
coordinates regular at the horizon.

The four-dimensional route below has two useful features. It needs only
two differentiated orders of asymptotic metric decay and one of tensor
decay. It also permits additional complete ends of the intermediate static
base until the classification argument excludes them. The required
positive-mass input is nonspin and asserts nonnegativity only; the
zero-mass rigidity and the treatment of the compactified point are proved
in Section~\ref{q4:cla:section}.

\subsection{Data, cuts, and the numerical input}

Let $\Omega$ be a smooth connected orientable four-manifold with nonempty
compact smooth boundary $S$. A smooth Riemannian metric $g$ is complete
with $S$ included, and a smooth symmetric covariant two-tensor $K$ is its
second fundamental form datum. We use the geometric density convention
\[
 \omega=\omega_3=|\mathbb S^3|=2\pi^2,\qquad
 \tau=\tr_gK,\qquad
 \mu=\tfrac12(R_g+\tau^2-|K|_g^2),\qquad
 J_i=\nabla^j(K_{ij}-\tau g_{ij}).
\]
Here $R_g=\Scal_g$ and $\Delta_g=\operatorname{div}_g\nabla$. The dominant energy
condition is $\mu\ge |J|_g$. The boundary normal $\nu$ points into the
exterior, and the future null expansion is
\[
 H=\operatorname{div}_S\nu,\qquad
 \theta_+=H+\tr_SK.
\]
For the second fundamental form in a spacetime, our sign convention is
\begin{equation}\label{q4:eq:intro-second-fundamental-form}
 K(U,V)=\overline g(\overline\nabla_U n,V)
\end{equation}
for the future unit normal $n$.

Assume that $\Omega$ has one end, diffeomorphic to the complement of a
closed ball in $\mathbb R^4$, with compact complement. In its coordinates
we require, for some $q>1$,
\begin{equation}\label{q4:eq:weak-decay}
 g-\delta=O_2(r^{-q}),\qquad K=O_1(r^{-1-q}),
 \qquad \mu,|J|_g\in L^1(\Omega,dV_g).
\end{equation}
The notation $O_j(r^{-a})$ includes coordinate derivatives of order
$k\le j$ with bounds $O(r^{-a-k})$. The ADM limits are assumed finite,
with normalization
\begin{align}
 E&=\frac1{6\omega}\lim_{R\to\infty}
 \int_{|x|=R}(\partial_jg_{ij}-\partial_i g_{jj})
              n_\delta^i\,dA_\delta,\label{q4:intro:energy}\\
 P_i&=\frac1{3\omega}\lim_{R\to\infty}
 \int_{|x|=R}(K_{ij}-\tau g_{ij})n_\delta^j\,dA_\delta.
 \label{q4:intro:momentum}
\end{align}

\begin{definition}[Full enclosing cuts]\label{q4:def:cuts}
A full enclosing cut is the entire intrinsic manifold boundary of a
connected smooth codimension-zero submanifold $D\subset\Omega$, closed
in $\Omega$, with compact boundary and with its manifold interior in
$\operatorname{int}\Omega$, that contains the whole sufficiently distant
end. Define
\[
 A_*(g)=\inf_D|\partial_{\mathrm{man}}D|_g.
\]
The cut may be disconnected, and every part coincident with $S$ is
counted. In particular $D=\Omega$ is allowed and has full cut $S$.
The infimum is not assumed attained by a smooth cut.
\end{definition}

The numerical theorem is used on a larger class than the equality theorem.
We record the exact one-ended input, so that no outermostness condition is
silently imposed on the variations.

\begin{theorem}[Weak four-dimensional numerical inequality]
\label{q4:thm:numerical-provider}
Suppose $(\Omega,g,K)$ satisfies the smoothness, orientability,
completeness, end, integrability, decay, and finite-charge assumptions
above, as well as $\mu\ge|J|_g$ and $\theta_+\le0$ on every component
of its nonempty compact boundary. Then $A_*(g)>0$ and
\begin{equation}\label{q4:eq:numerical}
 E\ge\sqrt{|P|_\delta^2+\frac14
                    \left(\frac{A_*(g)}\omega\right)^{4/3}}.
\end{equation}
In particular $E>|P|_\delta$ and
$\sqrt{E^2-|P|_\delta^2}\ge\frac12(A_*(g)/\omega)^{2/3}$.
\end{theorem}

This is the one-ended restriction of Theorem~8.2 in the numerical
companion \cite{CompanionNumerical}. The proof below never reverses an
end replacement or a graph deformation used to establish it.

\subsection{The equality class and its geometric conclusion}

\begin{definition}[Connected outermost horizon]\label{q4:def:horizon}
The data satisfy all hypotheses of
Theorem~\ref{q4:thm:numerical-provider}, have connected boundary $S$, and
satisfy the following additional conditions:
\begin{enumerate}[label=(\roman*),leftmargin=*]
\item $S$ is a future marginally outer trapped surface (MOTS):
      $\theta_+(S)=0$ for the normal into $\Omega$.
\item No smooth compact embedded two-sided enclosing hypersurface in
      $\operatorname{int}\Omega$, possibly disconnected and oriented toward
      the end, has $\theta_+\le0$ everywhere.
\item $S$ is outer area-minimizing: every full cut has area at least
      $A:=|S|_g>0$. Thus $A_*(g)=A$.
\end{enumerate}
\end{definition}

For $M_0>0$, the Schwarzschild--Tangherlini spacetime
\cite{Tangherlini1963} in its static exterior has metric
\begin{equation}\label{q4:eq:intro-tangherlini}
 \overline g_{M_0}
 =-\left(1-\frac{2M_0}{r^2}\right)dt^2
  +\left(1-\frac{2M_0}{r^2}\right)^{-1}dr^2
  +r^2g_{\mathbb S^3},\qquad r>\sqrt{2M_0}.
\end{equation}
Its regular maximal extension includes the future and past horizons and
their bifurcation sphere. The singular coefficient at $r=\sqrt{2M_0}$
in this display is removed by regular horizon coordinates.

\begin{theorem}[Original-data exterior rigidity]\label{q4:thm:rigidity}
Let $(\Omega,g,K)$ belong to Definition~\ref{q4:def:horizon}, and set
$m=\sqrt{E^2-|P|_\delta^2}$. Then
\begin{equation}\label{q4:eq:intro-equality}
 m=\frac12\left(\frac A\omega\right)^{2/3}
\end{equation}
if and only if the original initial data admit a global smooth spacelike
embedding, including $S$, into the regular maximal extension of
$\overline g_m$. The image lies in the closure of the corresponding
exterior, its end approaches that spatial infinity, and $S$ maps to a
smooth section of its future horizon or to its bifurcation sphere.
The induced metric is the original $g$, and its second fundamental form
for a consistent future unit normal is the original $K$ with convention
\eqref{q4:eq:intro-second-fundamental-form}. The embedding and normal are
smooth up to $S$ in regular horizon coordinates.

Equivalently, every such hypersurface in $\overline g_{M_0}$ satisfying
Definition~\ref{q4:def:horizon} and the prescribed asymptotic assumptions
has invariant ADM mass $M_0$, horizon area
$\omega(2M_0)^{3/2}$, and equality in \eqref{q4:eq:numerical}.
\end{theorem}

This is an exterior theorem: there is no hypothesis about an ambient
extension of the data behind $S$. Spherical topology, simple
connectivity, vacuum, and stationary development are conclusions of the
argument where they are used. The conclusion includes nonzero $K$ and
asymptotically boosted slices.

\subsection{How the original data are recovered}

The least enclosing area may have several minimizing frontiers.
Section~\ref{q4:enc:section} first constructs their compact space and
computes the right metric derivative as the smallest variation among
all of them. Positive separation in Section~\ref{q4:var:section} then
produces an area-sourced multiplier. Compact normal variations and
outermostness localize that source to $S$, before a homogeneous interior
adjoint equation is used.

Section~\ref{q4:adj:section} turns the multiplier into a smooth causal
lapse--shift field. Complementarity initially permits null stress;
vacuum is obtained only on its timelike region.
Section~\ref{q4:sta:section} verifies the hypotheses of the shared
static-base proposition at this weak decay and records the attachment
of the original horizon. The shared twist argument treats the possible
null set directly. Possible complete ends at interior zeros are retained
at this stage.

The conformal double in Section~\ref{q4:cla:section} follows the method of
Bunting and Masood-ul-Alam and its higher-dimensional development by
Gibbons, Ida, and Shiromizu
\cite{BuntingMasoodUlAlam1987,GibbonsIdaShiromizu2002}.
Second-order harmonic asymptotics make the compactification $C^2$.
A compact smoothing and conformal correction reduce zero-mass rigidity
to the smooth nonnegative-mass theorem of Lesourd, Unger, and Yau
\cite[Theorem~1.2 and Definition~1.9, arXiv version~1]{LesourdUngerYau2021},
with every other end still allowed. Euclidean rigidity then identifies
the entire original exterior. Finally
Section~\ref{q4:rec:section} reconstructs $g$ and $K$ as a single global
graph, makes the graph smooth at the horizon in the original boundary
structure, and computes the charges of every asymptotically boosted
slice for the converse.
